\documentclass[10pt,a4paper]{article}
\usepackage[margin=25mm]{geometry}
\usepackage{amsmath,amssymb,lmodern,booktabs}
\setlength{\emergencystretch}{2em}
\begin{document}
\begin{center}
{\Large A commutativity spot-check with detection probability $3/8$}\\[4pt]
{\large Counterexample to Conjecture 00000005449}\\[6pt]
ziangni-sys \qquad 4 October 2026
\end{center}

\section*{Claim and sampling law}
Both source versions assert that a single check has detection power
at least $1/2$ on a noncommutative witness. Consider the standard
commutativity spot-check: draw $X,Y$ independently and uniformly from
a finite group $G$, and detect noncommutativity if $XY\ne YX$.
Thus every ordered pair has probability $1/|G|^2$. We show that this
ordinary experiment has detection probability $3/8<1/2$ on a
genuine noncommutative group.

\section*{The actual group}
Let $G$ be the symmetry group of the square, of order eight, often
denoted $D_4$ or $D_8$ depending on the indexing convention.
Write its elements as
\[
\{r_i:i\in\mathbb Z/4\mathbb Z\}\cup
\{sr_i:i\in\mathbb Z/4\mathbb Z\}.
\]
With subscripts computed modulo four, multiplication is
\[
r_i r_j=r_{i+j},\quad r_i sr_j=sr_{j-i},\quad
sr_i r_j=sr_{i+j},\quad sr_i sr_j=r_{j-i}.
\]
In particular $r_1 sr_0=sr_3$ whereas $sr_0r_1=sr_1$,
so the group is noncommutative.

\section*{Complete pair count}
For each $x$, count its commuting partners using the displayed
multiplication laws:
\begin{center}
\begin{tabular}{lrrr}
\toprule
Elements $x$ & Number & Partners per element & Ordered pairs\\
\midrule
$r_0,r_2$ & 2 & 8 & 16\\
$r_1,r_3$ & 2 & 4 & 8\\
$sr_0,sr_1,sr_2,sr_3$ & 4 & 4 & 16\\
\bottomrule
\end{tabular}
\end{center}
All rotations commute with rotations. A rotation $r_i$ commutes with
a reflection precisely when $2i=0$ modulo four.
Two reflections $sr_i,sr_j$ commute precisely when $2(j-i)=0$,
so each reflection has two reflection partners and two rotation partners.
The table therefore covers all elements and gives $40$ commuting
ordered pairs out of $8^2=64$.
The actual detection event has $64-40=24$ pairs, so
\[
\mathbb P(XY\ne YX)=\frac{24}{64}=\frac38<\frac12.
\]
The possible equality $X=Y$ is included, as required for independent
uniform sampling; it is not removed by conditioning.
This refutes the single-check lower bound. The other clauses about
independent repetitions and automorphism corrections are not needed
for this disproof, and unspecified alternative algorithms are not claimed
to share this sampling law.

\section*{Formal verification}
Lean uses the actual Mathlib group \texttt{DihedralGroup 4}, with its
group structure and eight-element finite enumeration. It proves the
noncommuting witness and counts the full finite filters defined by
actual equality and inequality of the two products.
Each draw has rational mass $1/8$, and the pair mass is defined as
the product of the draw masses. Lean verifies nonnegative pair masses
and normalization, and defines detection probability by summing these
weights over the actual noncommuting-pair event.
The final theorem includes noncommutativity, a valid sampling
distribution, the exact probability, and failure of the proposed lower bound.
Finite checks use ordinary kernel-checked \texttt{decide};
no multiplication table or count is assumed and no native decision
shortcut or external numerical calculation is used.
\end{document}
